% Generated draft: requires mathematical and source review.
\documentclass[11pt]{article}
\usepackage[margin=25mm]{geometry}
\usepackage{fontspec}
\setmainfont{Latin Modern Roman}
\usepackage{amsmath,amssymb,mathtools,unicode-math}
\setmathfont{Latin Modern Math}
\usepackage{enumitem,needspace,xurl,hyperref}
\setlength{\parindent}{0pt}
\setlength{\parskip}{5pt}
\setlength{\emergencystretch}{3em}
\newcommand{\N}{\mathbb{N}}
\newcommand{\Z}{\mathbb{Z}}
\newcommand{\Q}{\mathbb{Q}}
\newcommand{\R}{\mathbb{R}}
\newcommand{\C}{\mathbb{C}}
\newcommand{\F}{\mathbb{F}}
\newcommand{\A}{\mathbb{A}}
\newcommand{\PP}{\mathbb{P}}
\newcommand{\E}{\mathbb{E}}
\newcommand{\eps}{\varepsilon}
\newcommand{\dd}{\,\mathrm{d}}
\begin{document}
\section*{UnsolvedMath ID: 701119}
\textbf{Source problem code:} catalogue entry
\subsection{Definitions and mathematical statement}
Let \(p\) be a prime and let \(G\) be a finite \(p\)-group. For \(g\in G\), define its conjugacy class and centralizer by
\[
g^G=\{x^{-1}gx:x\in G\},\qquad C_G(g)=\{x\in G:xg=gx\}.
\]
The breadth of \(g\) is the nonnegative integer
\[
b_G(g)=\log_p|g^G|=\log_p[G:C_G(g)].
\]
Define \(\gamma_1(G)=G\) and \(\gamma_{i+1}(G)=[\gamma_i(G),G]\), where \([H,G]\) is generated by the commutators \(h^{-1}x^{-1}hx\) with \(h\in H\) and \(x\in G\). Write \(c(G)\) for the nilpotency class, with \(c(\{1\})=0\); for nontrivial \(G\), this is the least positive integer \(c\) such that \(\gamma_{c+1}(G)=\{1\}\).

For a nonnegative integer \(n\), assume that there exists a generating set \(S\subseteq G\) such that
\[
G=\langle S\rangle,\qquad b_G(s)\le n\quad\text{for every }s\in S.
\]
Thus the breadth restriction applies to the chosen generators, not necessarily to every element of \(G\).

\textbf{Question.} Do there exist real constants \(A,B\ge0\), independent of \(p\), \(n\), and \(G\), such that every group satisfying these hypotheses obeys
\[
c(G)\le An+B?
\]
Here a linear bound is expressed as an affine upper bound. The supplied wording does not explicitly state whether the constants must be independent of the prime \(p\). The displayed question uses the uniform interpretation; under a prime-dependent interpretation, the quantifiers would instead be: for every prime \(p\), there exist \(A_p,B_p\ge0\) such that \(c(G)\le A_pn+B_p\) for every admissible \(n\) and finite \(p\)-group \(G\).

The supplied record attributes the bound \(c(G)\le n^2+1\) to M. R. Vaughan-Lee and J. Wiegold, and the improvement \(c(G)\le n^2-n+1\) to A. Mann. These are supplied background claims, not independently verified here.
\subsection{Short English statement}
If a finite \(p\)-group has a generating set whose individual generators each have at most \(p^n\) conjugates, must its nilpotency class admit a linear upper bound in \(n\)? The source does not explicitly specify whether the bound must be uniform over all primes. The registry labels this problem open; that label is supplied source data, not an independently verified assessment.
\subsection{Sources}
\begin{itemize}
\item Local UnsolvedMath registry, supplied source record with permanent ID \(701119\), exact source problem code \texttt{catalogue entry}, titled \emph{Kourovka Notebook Problem 15.31}. The supplied record labels the problem open.
\item Ya. G. Berkovich, proposer, 2002. \emph{The Kourovka Notebook (Unsolved Problems in Group Theory)}, 15th issue (2002), Problem 15.31; the supplied background also identifies the 21st issue edition (2026), \texttt{arXiv:1401.0300v48}. The supplied collection description names E. I. Khukhro and V. D. Mazurov as editors.
\item M. R. Vaughan-Lee and J. Wiegold, \emph{Proc. R. Soc. Edinburgh Sect. A} \textbf{95} (1983), 215--221. Cited by the supplied statement for the bound \(c(G)\le n^2+1\); no article title or link is supplied.
\item A. Mann, \emph{J. Group Theory} \textbf{4}, no. 3 (2001), 241--246. Cited by the supplied statement for the bound \(c(G)\le n^2-n+1\); no article title or link is supplied.
\end{itemize}
\end{document}
